% =====================================================================
% Section: Refinement-Typed Shape & Grad Calculus
% Target: NeurIPS 2026 main paper; column-figure-free; <= 2.5pp dense.
% Required preamble (in neurips.tex): amsmath, amssymb, amsthm, mathpartir
% (\inferrule), and the {theorem}/{definition}/{lemma} environments
% already present in docs/paper/neurips.tex.
% =====================================================================
\section{Refinement-Typed Shape \& Grad Calculus}
\label{sec:refinement-types}

We formalize TensorGuard's abstract domain as a refinement-type
system over a torch.fx-style intermediate representation
(IR)~\cite{reed2022torchfx}. The system has two novel ingredients
relative to prior shape-only refinement work
(\Cref{sec:related}): (i)~shape predicates live in a fragment of
linear integer arithmetic with bounded multiplication and
divisibility, decided by Z3~\cite{moura2008z3}; (ii)~each tensor
also carries a three-valued \emph{grad refinement} that tracks
whether the tensor is on the autograd tape, enabling sound
detection of silently frozen parameters.

\subsection{Types and judgments}
\label{sec:rt:types}

\paragraph{Tensor types.} The type of a tensor is
\[
  \mathsf{Tensor}[\tau,\;\varphi_{\mathsf{shape}}(d_1,\dots,d_k),\;\varphi_{\mathsf{grad}}],
\]
where $\tau\in\{\mathtt{f16},\mathtt{f32},\mathtt{f64},\mathtt{i32},\mathtt{i64},\mathtt{bool},\dots\}$,
the shape refinement
$\varphi_{\mathsf{shape}}\in\mathcal{P}_{\mathsf{shape}}$ is a
quantifier-free predicate over symbolic dim variables
$d_i\in\mathcal{V}$ in the theory
\[
  \mathcal{T}_{\mathsf{shape}} \;=\; \mathrm{LIA} \;\cup\;
  \mathrm{Div} \;\cup\; \mathrm{BMul},
\]
i.e.\ linear integer arithmetic, divisibility ($a\mid b$), and
bounded nonlinear multiplication (products whose factors are
syntactically known at extraction time, as enforced by the
canonicalization in
\texttt{src/typing\_rules.py}). The grad refinement
\(
  \varphi_{\mathsf{grad}}\in\{\mathsf{has\_grad},\,\mathsf{no\_grad},\,\top\}
\)
abstracts the conjunction
$\mathtt{requires\_grad}\wedge\mathtt{is\_on\_tape}$ where
$\mathtt{is\_on\_tape}$ is the runtime predicate ``some
differentiable operator producing this tensor was evaluated with
\texttt{torch.is\_grad\_enabled()=True}''.

\paragraph{Lattice.} The grad lattice is the flat lattice with
top element $\top$ and incomparable atoms
$\mathsf{has\_grad}$ and $\mathsf{no\_grad}$:
$\mathsf{has\_grad}\sqsubseteq\top$,
$\mathsf{no\_grad}\sqsubseteq\top$, and the two atoms are
\emph{not} ordered with respect to each other---they make
incompatible operational claims, so neither soundly substitutes
for the other. Joins are
$\mathsf{has\_grad}\sqcup\mathsf{no\_grad}=\top$. Subtyping on
tensor types is componentwise: $\tau$ and $k$ must match,
$\varphi_{\mathsf{shape}}$ is ordered by SMT entailment under the
ambient context, and $\varphi_{\mathsf{grad}}$ uses the lattice
above.

\paragraph{Contexts and entailment.} A typing context $\Gamma$
binds variables to refinement types and accumulates path
predicates and a grad-context flag $G\in\{\mathsf{on},\mathsf{off}\}$
recording whether the cursor is inside
\texttt{torch.no\_grad()}. The entailment judgment
$\Gamma\models\varphi\Rightarrow\psi$ is discharged by Z3 on the
shape projection of $\Gamma$; on the grad component it is the
lattice order. We write $\Gamma\vdash e:T$ for the typing
judgment.

\subsection{Typing rules for the supported fragment}
\label{sec:rt:rules}

We give representative rules; the full set (covering
\texttt{matmul, conv2d/3d, broadcast, view, reshape, split,
chunk, unbind, cat, stack, permute, transpose, einsum, indexing,
slicing, scatter, gather, repeat, expand,
F.scaled\_dot\_product\_attention, nn.LayerNorm, nn.RMSNorm}) is
implemented in \texttt{src/typing\_rules.py} and
\texttt{src/refinement/} (files \texttt{attention.py},
\texttt{norms.py}, \texttt{qkv.py}, \texttt{reshape.py}). Throughout,
we write $g(e_1,\dots,e_n) := \bigsqcup_i\varphi_{\mathsf{grad}}(e_i)$
when $G=\mathsf{on}$ and the operator is differentiable, and
$g(\dots):=\mathsf{no\_grad}$ otherwise; this captures PyTorch's
runtime semantics that any differentiable consumer of a tape
tensor under an active grad context produces a tape tensor.

{\small
\begin{mathpar}
\inferrule*[lab=T-Matmul]
{\Gamma\vdash a:\mathsf{Tensor}[\tau,\,\mathsf{shape}{=}(\bar B,m,k),\,\varphi_a]\\
 \Gamma\vdash b:\mathsf{Tensor}[\tau,\,\mathsf{shape}{=}(\bar B',k,n),\,\varphi_b]\\
 \Gamma\models\mathrm{broadcast}(\bar B,\bar B')=\bar C}
{\Gamma\vdash a\mathbin{@}b:\mathsf{Tensor}[\tau,\,(\bar C,m,n),\,g(a,b)]}
\and
\inferrule*[lab=T-Conv2d]
{\Gamma\vdash x:\mathsf{Tensor}[\tau,(B,C_{\mathrm{in}},H,W),\varphi_x]\\
 H'=\lfloor(H{+}2p_h{-}d_h(k_h{-}1){-}1)/s_h\rfloor{+}1\\
 W'=\lfloor(W{+}2p_w{-}d_w(k_w{-}1){-}1)/s_w\rfloor{+}1}
{\Gamma\vdash\mathrm{Conv2d}(x):\mathsf{Tensor}[\tau,(B,C_{\mathrm{out}},H',W'),g(x,\theta)]}
\and
\inferrule*[lab=T-Reshape-($-1$)]
{\Gamma\vdash x:\mathsf{Tensor}[\tau,(d_1,\dots,d_n),\varphi_x]\\
 \bar s=(s_1,\dots,s_{j-1},-1,s_{j+1},\dots,s_m)\\
 P=\textstyle\prod_i d_i\;\;Q=\textstyle\prod_{i\neq j}s_i\\
 \Gamma\models Q\mid P}
{\Gamma\vdash\mathrm{view}(x,\bar s):\mathsf{Tensor}[\tau,(s_1,\dots,P/Q,\dots,s_m),\varphi_x]}
\and
\inferrule*[lab=T-Cat]
{\forall i.\;\Gamma\vdash x_i:\mathsf{Tensor}[\tau,(d_1,\dots,d_k^{(i)},\dots,d_n),\varphi_i]}
{\Gamma\vdash\mathrm{cat}([x_i],\,\mathrm{dim}{=}k):
 \mathsf{Tensor}[\tau,(d_1,\dots,\textstyle\sum_i d_k^{(i)},\dots,d_n),\,\textstyle\bigsqcup_i\varphi_i]}
\and
\inferrule*[lab=T-Split]
{\Gamma\vdash x:\mathsf{Tensor}[\tau,(d_1,\dots,d_k,\dots,d_n),\varphi]\\
 \Gamma\models s\mid d_k\\\;\;m=d_k/s}
{\Gamma\vdash\mathrm{split}(x,s,k):
 (\underbrace{\mathsf{Tensor}[\tau,(\dots,s,\dots),\varphi],\dots}_{m\text{ copies}})}
\and
\inferrule*[lab=T-SDPA]
{\Gamma\vdash q,k:\mathsf{Tensor}[\tau,(B,H,T,D),\varphi_{q,k}]\\
 \Gamma\vdash v:\mathsf{Tensor}[\tau,(B,H,S,D),\varphi_v]\\
 \Gamma\models T=S}
{\Gamma\vdash\mathrm{SDPA}(q,k,v):\mathsf{Tensor}[\tau,(B,H,T,D),g(q,k,v)]}
\and
\inferrule*[lab=T-LayerNorm]
{\Gamma\vdash x:\mathsf{Tensor}[\tau,(d_1,\dots,d_n),\varphi_x]\\
 \mathrm{normalized\_shape}=(d_{n-r+1},\dots,d_n)}
{\Gamma\vdash\mathrm{LayerNorm}(x):\mathsf{Tensor}[\tau,(d_1,\dots,d_n),g(x,\gamma,\beta)]}
\and
\inferrule*[lab=T-Detach]
{\Gamma\vdash x:\mathsf{Tensor}[\tau,\varphi_s,\varphi_g]}
{\Gamma\vdash x.\mathrm{detach}():\mathsf{Tensor}[\tau,\varphi_s,\mathsf{no\_grad}]}
\and
\inferrule*[lab=T-NoGradCtx]
{\Gamma[G\mapsto\mathsf{off}]\vdash e:\mathsf{Tensor}[\tau,\varphi_s,\varphi_g]}
{\Gamma\vdash\mathtt{with~torch.no\_grad():}\;e:\mathsf{Tensor}[\tau,\varphi_s,\mathsf{no\_grad}]}
\end{mathpar}}

The remaining operators (\texttt{permute, transpose, einsum,
indexing, slicing, scatter, gather, repeat, expand, stack, chunk,
unbind, RMSNorm, conv3d}) follow the same template: shape is
computed by the corresponding handler in \texttt{src/refinement/}
or \texttt{src/typing\_rules.py}, and the grad component is
$g(\bar e)$ for differentiable ops and $\mathsf{no\_grad}$ for
non-differentiable indexing dtypes (e.g.\ integer indices in
\texttt{gather}). All rules degrade gracefully: if a precondition
is not entailed by $\Gamma$ but its negation is also not entailed,
the verdict is \textsc{Abstain} for that subterm and the result
type carries the symbolic shape unchanged.

\paragraph{Backward and optimizer rules.} The two
``whole-program'' rules close the loop:

{\small
\begin{mathpar}
\inferrule*[lab=T-Backward]
{\Gamma\vdash \mathit{loss}:\mathsf{Tensor}[\tau,(),\,\mathsf{has\_grad}]}
{\Gamma\vdash \mathit{loss}.\mathrm{backward}():\mathsf{unit}}
\and
\inferrule*[lab=T-OptStep]
{\forall p\in\mathit{opt}.\mathit{params}.\;\;
 \Gamma\vdash p:\mathsf{Tensor}[\tau_p,\varphi_p,\mathsf{has\_grad}]\\
 \text{i.e.\ } p\text{ is reachable from some }\mathit{loss}\text{ trace by a }\mathsf{has\_grad}\text{ path}}
{\Gamma\vdash \mathit{opt}.\mathrm{step}():\mathsf{unit}}
\end{mathpar}}

\textsc{T-Backward} requires the loss tensor to carry
$\mathsf{has\_grad}$; \textsc{T-OptStep} requires every parameter
the optimizer claims to update to be reachable from a
$\mathsf{has\_grad}$-typed expression in the trace. A parameter
typed $\mathsf{no\_grad}$ along every use in \texttt{forward} is
flagged: its \texttt{.grad} will be \texttt{None} after
\texttt{backward()} and the optimizer will silently skip it
(\Cref{sec:rt:gradflag}).

\subsection{Soundness}
\label{sec:rt:soundness}

The semantics is small-step, $\langle e,\sigma,G\rangle\to\langle
e',\sigma',G'\rangle$, where $\sigma$ is a heap of
$(\mathrm{shape},\mathrm{dtype},\mathtt{requires\_grad},\mathtt{is\_on\_tape})$
records and $G$ is the ambient grad context. We write
$\sigma\models\Gamma$ when every binding $x\!:\!T$ in $\Gamma$ has
$\sigma(x)$ satisfy the shape predicate of $T$ and the runtime
$(\mathtt{requires\_grad}\wedge\mathtt{is\_on\_tape})$ bit
implies $\varphi_{\mathsf{grad}}$ in the lattice sense
($\mathsf{true}$ may be typed $\mathsf{has\_grad}$ or $\top$;
$\mathsf{false}$ may be typed $\mathsf{no\_grad}$ or $\top$).

\begin{theorem}[Preservation up to subtyping]
\label{thm:preservation}
If $\Gamma\vdash e:T$, $\sigma\models\Gamma$, and $\langle
e,\sigma,G\rangle\to\langle e',\sigma',G'\rangle$, then there
exist $\Gamma'\supseteq\Gamma$ and $T'$ with
$\sigma'\models\Gamma'$, $\Gamma'\vdash e':T'$, and
$T'\sqsubseteq T$.
\end{theorem}
\begin{proof}[Proof sketch]
By induction on the reduction relation. For each non-context
rule, the operator's transfer function in
\texttt{src/typing\_rules.py} produces, by construction, a shape
that is the unique well-formed shape of the result up to SMT
equivalence; preservation on shapes is therefore
\emph{equality} modulo $\Gamma'$. For grad, the case
\textsc{T-NoGradCtx} discharges the obligation that the body
of \texttt{torch.no\_grad()} produces a tensor with
\texttt{is\_on\_tape}{=}false, so $\mathsf{no\_grad}$ is sound;
\textsc{T-Detach} likewise. For differentiable ops with $G=\mathsf{on}$,
PyTorch's runtime sets \texttt{is\_on\_tape}{=}true on the result
when at least one input has it true~\cite{paszke2019pytorch};
this is exactly the join $g(\bar e)$. Context rules use the IH
together with monotonicity of $g$ on the lattice. The relaxation
to $T'\sqsubseteq T$ is needed only at \textsc{T-NoGradCtx} exit
and at conditional joins where $\top$ is introduced.\qed
\end{proof}

\begin{theorem}[Three-valued progress]
\label{thm:progress}
Let $\mathsf{Verify}(e)\in\{\textsc{Verified},\textsc{Refuted},\textsc{Abstain}\}$
be TensorGuard's verdict. For any closed $e$ in the supported
fragment and any $\sigma\models\Gamma$:
(i) if $\mathsf{Verify}(e)=\textsc{Verified}$, then no reduction
sequence from $\langle e,\sigma,\mathsf{on}\rangle$ raises a
shape-mismatch exception; (ii) if
$\mathsf{Verify}(e)=\textsc{Refuted}$, then \emph{every} reduction
sequence raises a shape-mismatch exception at the witnessed
subterm; (iii) if $\mathsf{Verify}(e)=\textsc{Abstain}$, no claim
is made.
\end{theorem}
\begin{proof}[Proof sketch]
\textsc{Verified} means every operator's shape precondition is
entailed by $\Gamma$ via Z3; \Cref{thm:preservation} then
guarantees the precondition holds at the moment of evaluation, so
the runtime check inside the PyTorch dispatcher cannot fail.
\textsc{Refuted} requires Z3 to return a concrete model violating
some precondition; that model becomes the counterexample
$\sigma$, and the supported operators are
\emph{precondition-strict} (their PyTorch implementations raise
\texttt{RuntimeError} on the falsified condition).
\textsc{Abstain} corresponds to $\Gamma\not\models\varphi$ and
$\Gamma\not\models\neg\varphi$.\qed
\end{proof}

\begin{theorem}[Grad-flag soundness]
\label{thm:grad-soundness}
Suppose $\Gamma\vdash e:\mathsf{Tensor}[\tau,\varphi_s,\mathsf{has\_grad}]$
and $\sigma\models\Gamma$. Let $t$ be the tensor produced by
$\langle e,\sigma,\mathsf{on}\rangle\to^*\langle v,\sigma',\cdot\rangle$.
Then (a) $\mathrm{shape}(t)\models\varphi_s$ in $\sigma'$, and
(b) for every leaf parameter $p$ such that some directed path in
the autograd graph rooted at $t$ uses $p$ \emph{as a
differentiable input under} $G=\mathsf{on}$, executing
$t.\mathrm{backward}()$ leaves $p.\mathrm{grad}\neq\bot$.
\end{theorem}
\begin{proof}[Proof sketch]
(a) is the shape projection of \Cref{thm:preservation}. For (b),
$\varphi_{\mathsf{grad}}(t)=\mathsf{has\_grad}$ is, by induction
on the typing derivation and the rule
$g(\bar e)=\bigsqcup_i\varphi_{\mathsf{grad}}(e_i)$ (where the
join is taken only for differentiable ops with $G=\mathsf{on}$),
witnessed by a chain of $\mathsf{has\_grad}$ ancestors back to a
tape-leaf; PyTorch's autograd engine populates
$p.\mathrm{grad}$ for exactly those leaves
reachable by such a chain~\cite{paszke2019pytorch}. The
quantification over ``leaves used as differentiable inputs under
$G=\mathsf{on}$'' is essential: leaves used \emph{only} inside
\texttt{no\_grad} contexts have no such chain and are
\emph{correctly} predicted to receive no gradient by our system,
since their use sites type as $\mathsf{no\_grad}$.\qed
\end{proof}

\subsection{Grad-flag refinement in practice}
\label{sec:rt:gradflag}

Consider the model in
\texttt{experiments\_v5/grad\_flag\_example.py}: a
two-layer network whose trunk is computed inside
\texttt{with torch.no\_grad():} and whose head sits on top of
the (now-detached) trunk activation. By
\textsc{T-NoGradCtx}, the bound name $h$ types as
$\mathsf{Tensor}[\mathtt{f32},(16,8),\mathsf{no\_grad}]$. By
\textsc{T-Linear}, $\mathrm{head}(h)$ types as
$\mathsf{Tensor}[\mathtt{f32},(16,4),g(h,W_{\mathrm{head}})]$;
since $G=\mathsf{on}$ at the call site and
$\varphi_{\mathsf{grad}}(W_{\mathrm{head}})=\mathsf{has\_grad}$,
the join is $\mathsf{has\_grad}$. The loss therefore types as
$\mathsf{has\_grad}$ and \textsc{T-Backward} succeeds, but
$W_{\mathrm{trunk}}$ is reachable from the loss only via the
$\mathsf{no\_grad}$ binding $h$. The
\textsc{T-OptStep} side-condition
``$W_{\mathrm{trunk}}$ reachable by a $\mathsf{has\_grad}$ path''
is unsatisfiable, so TensorGuard reports
\textsc{Refuted} with the witness ``\texttt{trunk.weight.grad}
will be \texttt{None} after \texttt{backward()}.'' Running the
script reproduces the underlying PyTorch~2.9 behavior verbatim:
\texttt{head.weight.grad} is populated,
\texttt{trunk.weight.grad} is \texttt{None}, and
\texttt{optimizer.step()} leaves \texttt{trunk.weight}
bit-identical before and after the update.

\subsection{Limitations}
\label{sec:rt:limitations}

The calculus is calibrated to what we can prove. It does
\emph{not} cover: (i)~in-place mutation that re-routes the
autograd graph (\texttt{x.add\_(y)} when \texttt{y} carries grad),
which we conservatively type at $\top$; (ii)~user-defined
\texttt{torch.autograd.Function}, whose backward pass is opaque
and forces \textsc{Abstain}; (iii)~distributed and sharded
parameters (DDP, FSDP), where the ``unique path'' premise of
\Cref{thm:grad-soundness} fails because parameters are
materialized per-rank; (iv)~storage aliasing through
\texttt{view().detach()} chained with later in-place writes,
which can re-attach the storage to the tape in ways
$\varphi_{\mathsf{grad}}$ does not track; and
(v)~gradient-accumulation patterns across multiple
\texttt{backward()} calls without \texttt{zero\_grad()}, where
$.grad$ is non-\texttt{None} for reasons unrelated to the current
$e$. Cases (i)--(v) are detected at extraction time and downgrade
the relevant subterm to \textsc{Abstain}, preserving soundness of
\textsc{Verified}/\textsc{Refuted} verdicts at the cost of
completeness.

% =====================================================================
% Self-summary (for overseer agent):
% This section formalizes TensorGuard's abstract domain as a
% refinement-type system Tensor[tau, phi_shape, phi_grad] over a
% torch.fx-style IR. Shape refinements live in LIA + divisibility +
% bounded multiplication (Z3-decidable, matching the implementation in
% src/typing_rules.py and src/refinement/{attention,norms,qkv,reshape}.py).
% The grad refinement is a flat 3-valued lattice with has_grad and
% no_grad incomparable (both <: top); this ordering was corrected from
% an earlier draft after rubber-duck critique pointed out that the
% inverted ordering would let a no_grad tensor be used where has_grad
% is required, breaking soundness. Typing rules are given in mathpar
% style for matmul, conv2d, reshape (-1 inference via SMT divisibility),
% cat, split, SDPA, LayerNorm, detach, no_grad context, plus the two
% whole-program rules T-Backward and T-OptStep. We state Preservation
% up to subtyping (deliberately renamed from "Subject Reduction" per
% rubber-duck feedback, since shape preservation is exact but grad uses
% subtyping) and a 3-valued Progress theorem connecting Verified/
% Refuted/Abstain verdicts to the operational semantics. The grad-flag
% soundness theorem is tightened to quantify only over "leaves used as
% differentiable inputs under G=on" (per rubber-duck), avoiding the
% earlier overclaim. The worked example
% (experiments_v5/grad_flag_example.py) was actually executed under
% python3.11 + torch 2.9.1 and confirmed: head.weight.grad is
% populated, trunk.weight.grad is None, opt.step() leaves trunk
% bit-identical. Limitations explicitly enumerate the five edge cases
% the rubber-duck flagged: in-place mutation, custom autograd.Function,
% distributed/sharded params, view+detach storage aliasing, and grad
% accumulation across backward() calls without zero_grad. All
% operator references point at real files in src/. No claim is made
% about coverage outside the supported fragment.
% =====================================================================
